\paragraph{A Hilbert realization of the coupled finite action.}
The charged action also admits an interacting quantum Hamiltonian once a
quantization prescription is supplied. The kinetic energy depends on the
matter field and its gauge-dependent interpolation. Its reduction therefore
uses that kinetic metric, rather than the Maxwell mass matrix alone.
The construction retains the fixed cone, its unwrapped real edge
coefficients, the autonomous scalar action with \(m^2,g\geq0\), and
\(e\ne0\).

Write \(y=(a,\psi)\in\mathcal Y=\mathbb R^{42}\times
\mathbb C^{13}\), regarded as a real manifold. Let
\(W(a)\psi=\Psi_a\), \(S(a,\psi)b=\partial_a(W(a)\psi)[b]\),
and define
\begin{align}
 G_y(v,v)&=v_a^{\mathsf T}Mv_a+
              2\|W(a)v_\psi+S(a,\psi)v_a\|_{L^2(\Omega)}^2,
 \label{eq:whitney-interacting-metric}\\
 V(y)&=\frac12a^{\mathsf T}Ka+
 \|D_x\Psi_a\|_{L^2}^2+m^2\|\Psi_a\|_{L^2}^2
                         +\frac g2\|\Psi_a\|_{L^4}^4.
 \label{eq:whitney-interacting-potential}
\end{align}
The factor two in \(G\) follows from the normalization of the scalar
kinetic term. Nodal interpolation and \(M>0\) imply that \(G\)
is a smooth positive metric, including at \(\psi=0\).
The potential is smooth and nonnegative.

\begin{theorem}[Gauge reduction and a neutral interacting Hilbert space]
\label{thm:whitney-interacting-quantum}
The mean-zero real gauge group has a global smooth quotient
\[
 \mathcal Q=\ker(D^{\mathsf T}M)\times\mathbb C^{13}
          \cong\mathbb R^{30}\times\mathbb C^{13}.
\]
Eliminating the twelve mean-zero scalar-potential coefficients from
\eqref{eq:whitney-charged-action} defines a smooth positive metric
\(\gamma\) on \(\mathcal Q\). In fixed Euclidean-orthonormal slice
coordinates there are constants \(c_*,C_*>0\) such that
\begin{equation}
 \frac{c_*|w|^2}{1+|\psi|^2}\leq\gamma_y(w,w)
             \leq C_*(1+|\psi|^2)|w|^2.
 \label{eq:whitney-interacting-global-metric-bound}
\end{equation}
This metric is complete. With \(J_y=(0,i\psi)\)
and \(c\) the remaining constant scalar potential, the reduced action is
\begin{equation}
 L_{\mathrm{red}}(y,\dot y,c)
 =\frac12\gamma_y(\dot y+ecJ_y,\dot y+ecJ_y)-V(y).
 \label{eq:whitney-interacting-reduced-action}
\end{equation}
Its Hamiltonian at \(c=0\) is
\(H_{\mathrm{cl}}(y,p)=\tfrac12\gamma_y^{-1}(p,p)+V(y)\),
with the remaining Gauss constraint \(p(J_y)=0\).

Supply \(\hbar>0\), the Riemannian volume \(d\mu=d\mathrm{vol}_\gamma\),
and Laplace--Beltrami quantization. Then
\(\mathcal H=L^2(\mathcal Q,d\mu)\) has a nonnegative self-adjoint
Hamiltonian \(\widehat H\): the operator
\(-\hbar^2\Delta_\gamma/2+V\) on \(C_c^\infty(\mathcal Q)\)
is essentially self-adjoint, and its unique self-adjoint extension is its
Friedrichs realization.
The closed, nonzero subspace
\begin{equation}
 \mathcal H_0=\{f\in\mathcal H:
 f(a,e^{i\alpha}\psi)=f(a,\psi)\text{ in }L^2
                    \text{ for every }\alpha\in\mathbb R\}
 \label{eq:whitney-interacting-neutral-hilbert}
\end{equation}
reduces \(\widehat H\). Its restriction \(\widehat H_0\) is
self-adjoint on \(\mathcal D(\widehat H)\cap\mathcal H_0\)
and generates the strongly continuous unitary evolution
\(\exp(-it\widehat H_0/\hbar)\). Thus a single interacting finite
action supplies both a classical constrained Hamiltonian and, under the
declared prescription, a compatible neutral quantum state space.
\end{theorem}

\begin{proof}
Let \(\mathfrak g_0=\{\xi\in\mathbb R^{13}:\sum_i\xi_i=0\}\).
Its action is \((a,\psi)\mapsto(a+D\xi,e^{ie\xi}\psi)\).
Because the incidence matrix has kernel the constants,
\(D^{\mathsf T}MD\) is positive on \(\mathfrak g_0\).
Every orbit therefore meets \(\mathcal Q\) exactly once: solve
\(D^{\mathsf T}MD\xi=-D^{\mathsf T}Ma\) in
\(\mathfrak g_0\). This solution depends smoothly on \(a\).
Together with the inverse gauge action it gives a global product
description \(\mathcal Y\cong\mathcal Q\times\mathfrak g_0\),
so no singular quotient is needed for these twelve gauge directions.

For a gauge velocity define
\(R_y\xi=(D\xi,ie\xi\mathbin{\odot}\psi)\).
Differentiating the exact interpolation gauge identity gives
\begin{equation}
 S(a,\psi)D\xi+W(a)(ie\xi\mathbin{\odot}\psi)
                 =ie(W_0\xi)\Psi_a.
 \label{eq:whitney-interacting-vertical-identity}
\end{equation}
Consequently the full Lagrangian is exactly
\(L=G_y(\dot y+R_y\phi,\dot y+R_y\phi)/2-V(y)\).
This identity retains the derivative of the dressing; discarding it
would change both the scalar-potential equation and the reduced metric.
Time-independent gauge transformations are isometries of \(G\) and
preserve \(V\), by the same pointwise unit-phase covariance.

Choose any basis of \(\mathfrak g_0\). At \(y\in\mathcal Q\)
let \(R\) denote its vertical matrix and put
\(I_y=R^{\mathsf T}G_yR\). The matrix \(I_y\) is positive:
the Maxwell contribution gives
\(\xi^{\mathsf T}D^{\mathsf T}MD\xi>0\) for nonzero
\(\xi\in\mathfrak g_0\).
For \(w\in T_y\mathcal Q\), set
\begin{align*}
 \eta_*(w)&=-I_y^{-1}R^{\mathsf T}G_yw,\\
 \gamma_y(w,w)&=G_y(w,w)-
 (R^{\mathsf T}G_yw)^{\mathsf T}I_y^{-1}(R^{\mathsf T}G_yw).
\end{align*}
Completing the square yields the exact identity
\[
 G_y(w+R\eta,w+R\eta)
 =\gamma_y(w,w)+(\eta-\eta_*)^{\mathsf T}I_y(\eta-\eta_*).
\]
Thus \(\gamma\) is independent of the chosen vertical basis and
smooth. If \(\gamma_y(w,w)=0\), positivity of \(G\) gives
\(w=-R\eta_*\). Since \(w\) is tangent to the Coulomb slice,
\(D^{\mathsf T}MD\eta_*=0\), so \(\eta_*=0\) and
\(w=0\).

The stronger bound uses the actual dressed basis. On a tetrahedron \(T\),
choose a node maximizing \(|z_i|\). Where \(\lambda_i\geq3/4\),
unit-modulus phases and the triangle inequality give
\(|W(a)z|\geq(2\lambda_i-1)|z_i|\geq|z_i|/2\).
This region has volume \(|T|/64\), so
\[
 \|W(a)z\|_{L^2(T)}^2\geq\frac{|T|}{1024}|z_T|^2.
\]
Summing over tetrahedra gives \(\|W(a)z\|^2\geq c_W|z|^2\),
uniformly in all unwrapped \(a\), with
\(c_W=\min_T|T|/1024>0\). The fixed path coefficients also give
\(\|W(a)\|\leq C\) and \(\|S(a,\psi)\|\leq C|\psi|\).
Let \(T_c\) and \(B\) be orthonormal bases of the Coulomb and mean-zero
spaces, and write \(w=(T_cb,z)\). For every \(\eta\), the vertical
identity and \(T_c^{\mathsf T}MD=0\) give
\begin{align*}
 E_\eta&:=G(w+R_y(B\eta),w+R_y(B\eta))\\
 &=b^{\mathsf T}T_c^{\mathsf T}MT_cb+
   \eta^{\mathsf T}B^{\mathsf T}D^{\mathsf T}MDB\eta+2\|F_\eta\|^2,\\
 F_\eta&=Wz+ST_cb+ie(W_0B\eta)\Psi_a.
\end{align*}
Both displayed Maxwell matrices are positive. The remaining two terms in
\(F_\eta-Wz\) have norm at most
\(C|\psi|(|b|+|\eta|)\). Consequently
\(|b|^2+|z|^2\leq C(1+|\psi|^2)E_\eta\) for every \(\eta\).
Taking the minimum proves the lower bound in
\eqref{eq:whitney-interacting-global-metric-bound}; taking \(\eta=0\)
proves the upper bound. Constants concern this fixed mesh.
In global coordinates \(q\), the lower bound implies
\(\gamma\geq c_*(1+|q|^2)^{-1}I\). An escaping curve therefore has
length at least \(\sqrt{c_*}\) times its radial \(\operatorname{arsinh}|q|\)
variation and cannot have finite length. Smooth positivity on compact sets
then proves completeness.

Write \(\phi=\eta+c\mathbf1\), with \(\eta\in\mathfrak g_0\).
The twelve \(\eta\) equations are precisely the preceding minimization,
with \(w=\dot y+ecJ_y\). Substitution proves
\eqref{eq:whitney-interacting-reduced-action}. Elimination is variationally
valid because the eliminated derivative \(\partial L/\partial\eta\)
vanishes identically at its unique solution. The Legendre transform gives
\[
 p=\gamma_y(\dot y+ecJ_y,\,\cdot\,),\qquad
 H(y,p,c)=\frac12\gamma_y^{-1}(p,p)+V(y)-ec\,p(J_y).
\]
Variation of \(c\), with \(e\ne0\), gives the remaining constraint.
It is the constant component of the original thirteen Gauss equations,
not an extra neutrality assumption added after quantization.

For the analytic construction use the Hermitian form, initially on
\(C_c^\infty(\mathcal Q)\),
\begin{equation}
 q(f,k)=\frac{\hbar^2}{2}\int_{\mathcal Q}
       \gamma^{-1}(d\overline f,dk)\,d\mu
       +\int_{\mathcal Q}V\overline f k\,d\mu.
 \label{eq:whitney-interacting-quantum-form}
\end{equation}
Its domain is dense because \(\mu\) has a smooth positive density
in the global Euclidean coordinates. The form is nonnegative and
closable. Indeed, if \(f_n\to0\) in \(L^2\) and \(f_n\)
is form-Cauchy, its differentials and \(\sqrt V f_n\) have
respective \(L^2\) limits. Integration against compactly supported
smooth vector fields, using the metric divergence, shows that the first
limit is zero. Local boundedness of \(V\) shows that the second limit
is zero on every compact set and hence globally. Therefore
\(q(f_n,f_n)\to0\), proving closability.

Let \(\overline q\) be its closure. The closed-form representation
theorem gives a unique nonnegative self-adjoint operator for this
particular closed form; this is the Friedrichs realization of the
displayed differential operator. Explicitly, with the inner product
linear in its second argument,
\[
 \mathcal D(\widehat H)=\{f\in\mathcal D(\overline q):
 \exists b\in\mathcal H\ \forall k\in\mathcal D(\overline q),\
       \overline q(k,f)=\langle k,b\rangle\},\qquad
 \widehat Hf=b.
\]
The representation theorem and its symmetry version are standard
functional analysis \cite{IbortLledoPerezPardo2014}; the reduction and
coefficient hypotheses needed here have been verified above. This
argument identifies the Friedrichs realization. Completeness strengthens
this identification: apply the essential-self-adjointness theorem of
Shubin \cite[Theorem~1.1]{Shubin2001}, with zero magnetic one-form and
scalar potential \(2V/\hbar^2\). The manifold and measure are smooth,
the metric is complete, and this potential is smooth and nonnegative.
Thus the minimal operator has exactly one self-adjoint extension.

The residual circle \(T_\alpha(a,\psi)=(a,e^{i\alpha}\psi)\)
preserves the slice, its vertical spaces, \(G\), and \(V\). It
therefore preserves \(\gamma\), \(\mu\), and \(q\).
The pullbacks \(\mathcal U_\alpha f=f\circ T_{-\alpha}\)
are unitary and strongly continuous, first on compactly supported smooth
functions by dominated convergence and then on \(\mathcal H\) by
density. They preserve the closed form and its domain. Uniqueness in
the displayed operator-domain characterization gives
\(\mathcal U_\alpha\widehat H=\widehat H\mathcal U_\alpha\)
with domain preservation. Their bounded group average
\(P_0=(2\pi)^{-1}\int_0^{2\pi}\mathcal U_\alpha\,d\alpha\)
is the orthogonal projection onto \(\mathcal H_0\) and commutes
with the resolvent of \(\widehat H\). Hence this subspace reduces
the operator, and its restriction is self-adjoint. Smooth compactly
supported functions of \(a\) and \(\sum_i|\psi_i|^2\)
give nonzero invariant vectors; the subspace is infinite-dimensional.
The space \(P_0C_c^\infty(\mathcal Q)\) is an operator core
for the restriction: averaging preserves smoothness, its support lies
in the compact circle-saturation of the original support, and averaging
is contractive in the graph norm because it commutes with \(\widehat H\).
Apply \(P_0\) to the compactly supported smooth graph approximants supplied
by essential self-adjointness. It is also a form core.
Spectral calculus supplies the stated strongly continuous unitary group.
For initial states in \(\mathcal D(\widehat H_0)\) it solves
\(i\hbar\partial_t f=\widehat H_0f\) in the strong sense; for
every Hilbert state it supplies norm-preserving evolution.
\end{proof}

With the convention
\(\mathcal U_{es}=\exp(-is\widehat Q/\hbar)\), the
self-adjoint charge generator acts
on compactly supported smooth functions as
\[
 \widehat Q=-i\hbar e\sum_i
       \left(u_i\frac{\partial}{\partial v_i}
                  -v_i\frac{\partial}{\partial u_i}\right),
 \qquad \psi_i=u_i+iv_i.
\]
Thus \(\mathcal H_0=\ker\widehat Q\) expresses the quantum
constant Gauss constraint. It allows matter amplitudes and neutral
correlations; it does not freeze the scalar field at zero.
The point \(\psi=0\), where the circle has a stabilizer, remains
an ordinary point of \(\mathcal Q\). Taking circle-invariant Hilbert
vectors avoids imposing a smooth structure on a singular full orbit space.

\begin{proposition}[A normalized neutral initial state]
\label{prop:whitney-interacting-gaussian-state}
In the orthonormal \(56\)-dimensional slice coordinates write
\(d\mu=\rho(q)\,dq\), with \(\rho=\sqrt{\det\gamma}\).
For any supplied \(\sigma>0\), the function
\begin{equation}
 f_\sigma(q)=\rho(q)^{-1/2}(2\pi\sigma^2)^{-14}
                    \exp\!\left(-\frac{|q|^2}{4\sigma^2}\right)
 \label{eq:whitney-interacting-gaussian-state}
\end{equation}
is a normalized vector in \(\mathcal D(\widehat H_0)\). Its initial
multiplication observables satisfy
\begin{align}
 \mathbb E\|\Psi_a\|_{L^2}^2&=\frac45\sigma^2|\Omega|,&
 \mathbb E\|\Psi_a\|_{L^4}^4&=\frac{48}{35}\sigma^4|\Omega|,
 \label{eq:whitney-interacting-gaussian-matter-moments}\\
 \mathbb E\frac{a^{\mathsf T}Ka}{2}
       &=\frac{\sigma^2}{2}\operatorname{tr}(T_c^{\mathsf T}KT_c).
 \label{eq:whitney-interacting-gaussian-magnetic-moment}
\end{align}
\end{proposition}
\begin{proof}
The measure \(|f_\sigma|^2d\mu\) is exactly
\(N(0,\sigma^2I_{56})\). The residual circle acts orthogonally and
preserves \(\rho\), so \(f_\sigma\) is normalized and neutral.

For the operator-domain condition, every derivative of \(W\) in \(a\)
is uniformly bounded, and the corresponding derivatives of \(S\) are
linear in \(\psi\). The vertical Gram matrix has the fixed positive
lower bound \(B^{\mathsf T}D^{\mathsf T}MDB\). Differentiating its inverse
and the Schur formula shows polynomial growth of all metric derivatives.
The global metric bound gives polynomial growth of \(\gamma^{-1}\)
and its derivatives; \(\partial_j\log\rho=
\tfrac12\operatorname{tr}(\gamma^{-1}\partial_j\gamma)\) gives the same
control for derivatives of \(\log\rho\). The potential and its
derivatives grow polynomially, with \(V(q)\leq C(1+|q|)^4\).
Thus \(f_\sigma\) and the differential expression
\((-\hbar^2\Delta_\gamma/2+V)f_\sigma\) lie in \(L^2(d\mu)\):
after multiplication by \(\sqrt\rho\), each differentiated term is
bounded by a polynomial times a Gaussian. Integration by parts against
\(C_c^\infty\) puts the state in the adjoint domain; essential
self-adjointness identifies that domain with \(\mathcal D(\widehat H)\).

Conditional on \(a\), independent circular Gaussian nodal coefficients
give pointwise second and fourth moments
\(2\sigma^2\sum_i\lambda_i^2\) and
\(8\sigma^4(\sum_i\lambda_i^2)^2\). The simplex averages of
\(\sum_i\lambda_i^2\) and \((\sum_i\lambda_i^2)^2\) are
\(2/5\) and \(6/35\), respectively. This proves the matter identities;
the magnetic identity follows from the transverse Gaussian covariance.
\end{proof}
This is a selected initial state, without a ground-state,
physical-preparation, or computed-time-history assertion.

The metric coefficients and their Schur reduction are approximated by
positive element quadrature in
\path{code/electromagnetism/whitney_interacting_quantum.py}; independent
tests check the actual cone, gauge covariance, scalar normalization,
Gauss elimination, and the nondegenerate zero-matter limit. These finite
checks compare exact simplex moments at \(a=0\) and quadrature orders
four, five and six at a fixed nonzero configuration. The latter comparison
measures quadrature sensitivity, not a certified integration error. The
tests do not prove closability or self-adjointness; those are the analytic
arguments above. The quantum prescription chooses a measure, operator
ordering and \(\hbar\); the extension is unique for that operator.
The code evaluates Gaussian half-densities and exact moment formulae;
its curved-measure amplitudes use the approximate metric density.
No equivalence with quantization before reduction, unique quantization,
quantum trajectory computation, continuum interacting field theory,
physical-state selection, or Born-rule selection follows.
The established attachment is to the same finite
charged action and its full kinetic metric.
